%%%%%%%%%%%%%%%%%%%%%%%%%%%%%%%%%%%%%%%%%%%%%%%%%%%%%%%%%%%%%%%%%%%%%%%%%%%%%%%
%% Estilo do modelo de monografia do iCEV
%%
%% Transcrito de monografia-tcc/main.tex: fontes de título, logo no cabeçalho,
%% capa, folha de rosto, ambiente de resumo, citação nas referências e cores
%% dos links. Divergências deliberadas do modelo, todas pelas normas ABNT:
%%   - Referências sem indicativo numérico (o modelo as numerava como capítulo).
%%   - Texto preto, inclusive nos links (NBR 14724, 5.1); o modelo usa azul.
%%   - Títulos em corpo 12, graduados por caixa-alta e negrito, com o mesmo destaque
%%     no sumário e no texto (NBR 14724, 5.1 e 5.4; NBR 6027). Por isso as seções
%%     secundárias não usam caixa-alta (sem section=TITLE em main.tex).
%%   - Um espaço entre indicativo e título (NBR 14724, 5.2.2).
%%   - Margem inferior de 2 cm e número da página a 2 cm da borda superior (NBR 14724,
%%     5.1 e 5.3); o modelo chega a cerca de 1,5 cm e a 2,4 cm.
%%   - Introdução em caixa-alta no sumário, como os demais capítulos.
%%   - Sem título corrente no cabeçalho: o logo de 15 cm cobre quase toda a
%%     largura do texto (16 cm) e escondia o título atrás da imagem. A NBR 14724
%%     só exige a paginação no canto superior direito.
%%   - Introdução numerada, como seção primária (NBR 14724, 5.2.3; NBR 6024); o modelo a
%%     deixa sem número. Decisão do usuário, em 2026-11-01.
%%   - Referências alinhadas à esquerda, com título em negrito (NBR 6023).
%% Os dados (título, autor, curso) ficam em preambulo/dados.tex.
%%%%%%%%%%%%%%%%%%%%%%%%%%%%%%%%%%%%%%%%%%%%%%%%%%%%%%%%%%%%%%%%%%%%%%%%%%%%%%%

%% --- Títulos -----------------------------------------------------------------
\renewcommand{\ABNTEXchapterfont}{\rmfamily\bfseries}
\renewcommand{\ABNTEXchapterfontsize}{\normalsize}
\renewcommand{\ABNTEXsectionfontsize}{\normalsize}
\renewcommand{\ABNTEXsubsectionfont}{\rmfamily\mdseries}
\renewcommand{\ABNTEXsubsectionfontsize}{\normalsize}
\setsecnumformat{\csname the#1\endcsname\ }
%% Um espaço entre número e título (NBR 14724, 5.2.2): a definição da classe deixa dois.
\makeatletter
\renewcommand{\printchapternum}{%
	\tocprintchapter
	\setboolean{abntex@innonumchapter}{false}%
	\chapnumfont
	\ifthenelse{\boolean{abntex@apendiceousecao}}%
		{\space\thechapter\tocinnonumchapter\ABNTEXcaptiondelim}%
		{\thechapter\space}%
}
\makeatother
%% Título separado do texto por uma linha em branco de espaço 1,5 (NBR 14724, 5.2.2).
\setlength{\afterchapskip}{0pt}
%% Destaque do sumário igual ao do texto: terciária sem negrito (NBR 14724, 5.4).
\renewcommand{\cftsubsectionfont}{\normalfont\normalsize}
\renewcommand{\cftsubsectionpagefont}{\normalfont\normalsize}
%% Apêndice no sumário: "APÊNDICE A --- TÍTULO".
\cftinsertcode{A}{%
	\renewcommand{\cftchapteraftersnum}{~\textemdash~}%
	\renewcommand{\cftchapterpresnum}{\unskip}%% a classe deixa dois espaços após "APÊNDICE"
}
\renewcommand{\tocinnonumchapter}{%
	%% 2 em cobre "A ---"; o memoir soma a largura de "APÊNDICE", e a segunda linha alinha sob
	%% o início do título.
	\addtocontents{toc}{\cftsetindents{chapter}{\cftlastnumwidth}{2em}}%
	\cftinserthook{toc}{A}%
}
%% O sumário já recebe o espaço que antecede a primeira entrada.
\renewcommand{\aftertoctitle}{\par\vspace{6.3pt}}
\renewcommand{\afterloftitle}{\par\vspace{\baselineskip}}
\renewcommand{\afterlottitle}{\par\vspace{\baselineskip}}
\setbeforesecskip{\baselineskip}
\setaftersecskip{\baselineskip}
\setbeforesubsecskip{\baselineskip}
\setaftersubsecskip{\baselineskip}
%% Equações em numeração corrida, como figuras e tabelas.
\counterwithout{equation}{chapter}

%% --- Notas, ilustrações e tabelas ---------------------------------------------
%% Filete de 5 cm a partir da margem esquerda (NBR 14724, 5.2).
\renewcommand{\footnoterule}{\kern-3pt\hrule width 5cm\kern 2.6pt}
%% Tabela em corpo 12; só fonte, legenda e nota ficam menores (NBR 14724, 5.1).
\renewcommand{\IBGEtabfontsize}{\normalsize}
%% Com o título acima, o caption não deixa espaço antes da fonte.
\LetLtxMacro{\fonteabntex}{\fonte}
\renewcommand{\fonte}[2][\fontename]{\par\vspace{3pt}\fonteabntex[#1]{#2}}
%% Legenda abaixo da fonte (NBR 14724, 5.8); o título vem antes da ilustração.
\newcommand{\legenda}[1]{\nota[Legenda]{#1}}
%% Mesma autoria, datas diferentes: (Autor, 2009, 2022), como pede a NBR 10520, 6.1.
\newcommand{\citemesmoautor}[3]{(#1, \citeyear{#2}, \citeyear{#3})}

%% --- Referências -------------------------------------------------------------
%% Sem "Citado na página" do modelo: a NBR 6023 não prevê esse elemento.
\renewcommand{\bibsection}{%
	\chapter*{\bibname}%
	\addcontentsline{toc}{chapter}{\MakeTextUppercase{\bibname}}%
	\bibmark
	\vspace*{21.4pt}%% uma linha em branco de espaço 1,5 (NBR 14724, 5.2.2)
	\prebibhook
}

%% --- Cabeçalho com logo --------------------------------------------------------
\newcommand{\logoicevheader}{\includegraphics[width=15cm]{imagens/logo-icev.jpg}}

\makepagestyle{icevpretextual}
\makeevenhead{icevpretextual}{}{\logoicevheader}{}
\makeoddhead{icevpretextual}{}{\logoicevheader}{}
\aliaspagestyle{titlingpage}{icevpretextual}

\makepagestyle{icevheadings}
%% O número sobe para ficar a 2 cm da borda superior (NBR 14724, 5.3).
\newcommand{\icevnumeropagina}{\raisebox{0.35cm}{\ABNTEXfontereduzida\thepage}}
\makeevenhead{icevheadings}{\icevnumeropagina}{\logoicevheader}{}
\makeoddhead{icevheadings}{}{\logoicevheader}{\icevnumeropagina}
\makeheadrule{icevheadings}{\textwidth}{\normalrulethickness}

\makepagestyle{icevchapfirst}
\makeoddhead{icevchapfirst}{}{\logoicevheader}{\icevnumeropagina}

%% O modelo usa 2.1cm, mas o logo mede 60,6pt e não cabe (aviso do memoir e
%% \vbox overfull em toda página). Com 2.15cm cabe, e cabeçalho mais
%% \headsep continuam dentro da margem superior de 3 cm.
\setlength{\headheight}{2.15cm}
\setlength{\headsep}{0.5cm}
\checkandfixthelayout
%% Sem este ajuste, a última linha de uma página cheia fica a cerca de 1,7 cm da borda;
%% a margem inferior é de 2 cm (NBR 14724, 5.1).
\addtolength{\textheight}{-0.5cm}
\checkandfixthelayout[fixed]

%% --- Links ---------------------------------------------------------------------
%% pdftitle e pdfauthor ficam em dados.tex: aqui \@title e \@author ainda estão vazios.
\makeatletter
\hypersetup{
	pdfsubject  = {Trabalho de Conclusão de Curso},
	pdfcreator  = {LaTeX with abnTeX2},
	pdfkeywords = {manutenção preditiva; análise de vibração; detecção de anomalias;
	               sistemas embarcados; instrumentação de baixo custo},
	hidelinks,
	bookmarksdepth = 4
}
\makeatother

%% --- Parágrafo -----------------------------------------------------------------
\setlength{\parindent}{1.25cm}
\setlength{\parskip}{0pt}

%% --- Capa ----------------------------------------------------------------------
\renewcommand{\imprimircapa}{%
	\begin{capa}
		\centering
		{\ABNTEXchapterfont\bfseries\instituicaoicev\par}
		{\ABNTEXchapterfont\bfseries\cursoicev\par}

		\vfill
		{\ABNTEXchapterfont\bfseries\imprimirautor}
		\vfill

		\begin{center}
			\ABNTEXchapterfont\bfseries\imprimirtitulo
		\end{center}
		\vfill

		{\bfseries\imprimirlocal\par}
		{\bfseries\imprimirdata}
		\vspace*{1cm}
	\end{capa}
}

%% --- Natureza do trabalho --------------------------------------------------------
%% Do meio da mancha até a margem direita, em espaço simples, na folha de rosto e na
%% de aprovação (NBR 14724, 5.2).
\newcommand{\icevnatureza}{%
	{\raggedleft
	\begin{minipage}{.5\textwidth}
		\SingleSpacing
		%% Espaços sem encolher: nesta coluna estreita, a linha apertada colava
		%% as palavras ao copiar do PDF. Sem hifenizar o nome da instituição.
		\spaceskip=\fontdimen2\font plus \fontdimen2\font\relax
		\hyphenpenalty=10000
		\imprimirpreambulo
	\end{minipage}\par}%
}

%% --- Folha de rosto ------------------------------------------------------------
\makeatletter
\renewcommand{\folhaderostocontent}{
	\begin{center}

		{\ABNTEXchapterfont\bfseries\imprimirautor}

		\vspace*{\fill}\vspace*{\fill}
		\begin{center}
			\ABNTEXchapterfont\bfseries\imprimirtitulo
		\end{center}
		\vspace*{\fill}

		\abntex@ifnotempty{\imprimirpreambulo}{%
			\icevnatureza
			\vspace*{\fill}
		}%

		{\abntex@ifnotempty{\imprimirinstituicao}{\imprimirinstituicao\vspace*{\fill}}}

		{\imprimirorientadorRotulo~\imprimirorientador\par}
		\abntex@ifnotempty{\imprimircoorientador}{%
			{\imprimircoorientadorRotulo~\imprimircoorientador}%
		}%
		\vspace*{\fill}

		{\bfseries\imprimirlocal}
		\par
		{\bfseries\imprimirdata}
		\vspace*{1cm}

	\end{center}
}
\makeatother

%% --- Resumo --------------------------------------------------------------------
\renewenvironment{resumo}[1][\resumoname]{%
	\PRIVATEbookmarkthis{#1}
	\renewcommand{\abstractnamefont}{\ABNTEXchapterfont\bfseries}
	\renewcommand{\abstractname}{\ABNTEXchapterupperifneeded{#1}}
	\begin{abstract}
}{\end{abstract}\PRIVATEclearpageifneeded}
\setlength{\absparsep}{18pt}
\addtolength{\abstitleskip}{1.7pt}

%% --- Trechos que mudam com as medições do dispositivo --------------------------
%% Definido aqui porque o resumo e o abstract também usam, antes das seções.
\providecommand{\emandamento}[1]{\textcolor{red}{#1}}

%% --- Quadros --------------------------------------------------------------------
%% Quadro é ilustração (NBR 14724, 5.8), com lista própria antes da de tabelas.
\newcommand{\quadroname}{Quadro}
\newcommand{\listofquadrosname}{Lista de quadros}
\newfloat[chapter]{quadro}{loq}{\quadroname}
\newlistof{listofquadros}{loq}{\listofquadrosname}
\newlistentry{quadro}{loq}{0}
\counterwithout{quadro}{chapter}
%% Mesmo separador das tabelas e figuras do abnTeX2 (espaço, meia-risca, espaço).
\captionsetup[quadro]{position=top,labelsep=endash}
%% Separador das legendas nos apêndices, onde \ABNTEXcaptiondelim vira travessão para o título
%% "APÊNDICE A --- TÍTULO"; \kxdelimlegenda guarda o separador do texto (main.tex, apendicesenv).
\DeclareCaptionLabelSeparator{kxlegenda}{\kxdelimlegenda}
\renewcommand{\cftquadroname}{\quadroname\space}
\renewcommand*{\cftquadroaftersnum}{\hfill\textendash\hfill}
\renewcommand{\afterloqtitle}{\par\vspace{\baselineskip}}

%% --- Referências -------------------------------------------------------------
%% Separadas por uma linha em branco de espaço simples (NBR 6023, 6.3).
\setlength{\bibitemsep}{10pt}
%% Data provável, entre colchetes (NBR 6023, 8.6): os colchetes no campo year quebram o
%% argumento opcional de \bibitem, por isso ela fica em uma macro.
\newcommand{\anomafaulda}{[201-?]}

%% --- Listas de siglas e de símbolos ------------------------------------------
%% Mesma separação entre título e texto das demais listas.
\renewcommand*{\glossarypreamble}{\vspace*{\baselineskip}}
\renewcommand{\nompreamble}{\vspace*{20pt}}

%% --- Páginas só de ilustrações -------------------------------------------------
%% Uma página com uma ou mais ilustrações e sem texto começa no topo, em vez de deixar
%% a ilustração centralizada no meio da página com espaço vazio acima.
\makeatletter
\setlength{\@fptop}{0pt}
\setlength{\@fpsep}{\floatsep}
\setlength{\@fpbot}{0pt plus 1fil}
\makeatother
